\documentclass{article}

\usepackage{amsfonts}

\begin{document}

Andrew Fowler • Mark McGuinness
Chaos
An Introduction for Applied Mathematicians

\hfill

Chapter 1
Introduction

\hfill

\textbf{1.1 Limit Cycles}


\hfill

The behaviour of autonomous (that is, where time is not represented explicitly)
ordinary differential equations in $\mathbb{R}^2$ is well understood. 

\hfill

The qualitative behaviour
of the equation

\[
\dot{x} = f (x), x \in \mathbb{R}^2
\]

is organised by its \textit{fixed points}, where $\dot{x} = f (x) = 0$. 

\hfill

Elementary texts show how
these maybe classified into \textit{saddles, spirals and nodes} depending on the eigenvalues
of the Jacobian Df(x) at the fixed point. 

\hfill

In particular, if these eigenvalues are complex
conjugates, there is a spiral, or focus; if they are real and of opposite sign, there
is a saddle; and if they are real and of the same sign, there is a node: see Fig. 1.1.

\hfill

Nodes and spirals may be stable or unstable depending on whether the eigenvalues $\lambda$
have $Re \lambda < 0$ or $Re \lambda > 0$, respectively; saddles are always unstable. 

\hfill

Other degenerate
cases (centre, degenerate node) can occur, if $Re \lambda = 0$, or if two eigenvalues
coalesce.

Knowledge of the behaviour of trajectories near fixed points informs the method
of phase plane analysis. 

\hfill

An example of its use is to find travelling solitary wave
solutions of the Korteweg–de Vries equation

\[
u_t + uu_x + u_{xxx} = 0
\]

in the form $u = f (x - ct) = f (\xi)$, so that

\[
- cf' + f f' + f''' = 0
\]

with first integral

\[
- cf + \frac{1}{2} f^2 + f'' =0 
\]

(so that $f(\pm \infty) = 0)$. 

\hfill

This represents a nonlinear oscillator with potential $V( f ) = - \frac{1}{2} cf^2 + \frac{1}{6} 
f^3$.1 

\hfill

If we take $c > 0$ (without loss of generality), then the phase plane is
as in Fig. 1.2. 

\hfill

There are two fixed points, f = 0, and f = 2c, which are respectively
a saddle and a centre. (The latter persists as a centre in the presence of nonlinear
terms by virtue of the existence of a first integral, which in turn is associated with the
absence of a damping term.) 

\hfill

Solitary waves exist, by virtue of the homoclinic orbit
which connects 0 to itself. 

\hfill

We see in addition that there exists a family of periodic
orbits lying within the homoclinic loop. One can think of the homoclinic orbit as a
limiting member of this family with infinite period.

More generally, periodic behaviour occurs in two-dimensional systems at isolated
limit cycles. A limit cycle is a closed curve in a phase space (which, therefore,
represents a periodic time series) to which nearby trajectories tend either as $t \to + \infty$
(stable limit cycle) or as $t \to \infty$ (unstable limit cycle). A simple example (see also
exercise 1.7) is offered by the system

\[
\dot{x} = x - y - x(x^2 + y^2),
\]
\[
\dot{y} = x + y - y(x^2 + y^2); 
\]


by transforming to polar coordinates $(r, \theta),$ it is straightforward to find the phase
plane shown in Fig. 1.3. There is a stable limit cycle at r = 1, and an unstable focus
at r = 0.
One is often interested in the behaviour of a (physical) system as a parameter
changes (e.g. the Earth’s climate as the atmospheric CO2 level changes). In this case,
one might study differential equations of the form

\[
\dot{x} = f (x, \mu), (1.6)
\]

where $\mu$ is a parameter. For some ranges of $\mu$, the system may exhibit steady
behaviour, in others it may behave periodically. One then wishes to chart regions
in parameter space where such qualitative changes occur. These delineate particular
values of $\mu$ at which such bifurcations take place, and $\mu$ is then called a bifurcation
parameter.

1.2 Bifurcation Theory

The Watt governor (see Fig. 1.4) is a device invented by James Watt around 1788 to
control the flow of steam into the cylinders of an engine. 

It is a simple but efficient
device, which acts on the basis of bifurcation theory. 

As shown in the figure, the
device consists in essence of a pair of rigid pendulums attached to a central shaft,
which rotates about its axis at a rate determined by the pressure of steam in the engine.


The object is to have a device such that if this rotation rate becomes too large, then
the arms of the device rise, thus opening a valve which alleviates the pressure.


If we denote the angle of the arms of the governor to the downward vertical as $\theta$,
and the angular velocity of the shaft as $\Omega$, then the equation governing the dynamics
is, ignoring friction (see question 5.1),

\[
\ddot{\theta} - \Omega^2 \sin \theta \cos \theta + \Omega^2_c \sin \theta = 0, (1.7)
\]


where

\[
\Omega_c = \sqrt{\frac{g}{l}}
\]

g is the acceleration due to gravity and l is the length of the pendulum arms.

The steady states of the system are $\theta = 0, \pi$, and also $\pm \cos^{-1} (\frac{\Omega^2_c}{\Omega^2})$ if $\Omega > \Omega_c$.

The function of the governor can be seen in the fact that we can have $theta$ as an
increasing function of $\Omega$ if $\Omega > \Omega_c$, thus opening the valve, and its operation relies
on the existence of a pitchfork bifurcation at $\Omega = \Omega_c$. More precisely, if we linearise
(1.7) near the origin, we see that

\[
\ddot{\theta} \approx (\Omega^2 - \Omega^2_c) \theta
\]

and thus the vertical state loses stability for $\Omega > \Omega_c$.

Expanding further in $theta$ and also
supposing that $\Omega - \Omega_c$ is small, we have

\[
\ddot{\theta} \approx (\Omega^2 - \Omega^2_c) \theta - \frac{1}{2} \Omega^2_c \theta^3...,
\]

which shows explicitly that there is an exchange of stability between $\theta = 0$ and the
non-zero fixed points. Bifurcation theory is discussed in more detail in Chap. 3.


\end{document}

